\section{Background and Related Work}
\label{sec:background}

\subsection{Searchable Encryption and Leakage}

Searchable symmetric encryption allows a client to query an encrypted index. Its security is
stated relative to a leakage function that describes what the scheme is permitted to reveal,
typically the access pattern (which encrypted records a query touches) and the search pattern
(whether two queries are equal) \cite{curtmola2006searchable,kamara2018structured}. A large body
of work shows that such leakage can be exploited. Islam, Kuzu, and Kantarcioglu recovered queries
from access-pattern co-occurrence \cite{islam2012access}; Cash et al. introduced the count attack
\cite{cash2015leakage}; and Blackstone, Kamara, and Moataz reduced the auxiliary knowledge these
attacks require \cite{blackstone2020revisiting}. Volume leakage alone suffices for database
reconstruction \cite{kellaris2016generic,grubbs2018pump}, and query recovery can survive
access-pattern obfuscation \cite{oya2022ihop}. Kornaropoulos, Papamanthou, and Tamassia recover
data from the leakage profile of $k$-nearest-neighbor queries \cite{kornaropoulos2019knn}, the
query type an encrypted vector store serves. Volume-hiding constructions pad response sizes to
remove this channel \cite{kamara2019volume,patel2019mitigating}. Xu et al. showed that
leakage-abuse attacks remain effective against schemes that are both forward and backward private
\cite{xu2023leakage}.

\subsection{Forward Privacy, Backward Privacy, and Injection}

Forward privacy requires that an update cannot be linked to previously issued search tokens
\cite{stefanov2014practical,bost2016sophos}. It was motivated by file-injection attacks, in which
an adversary who can insert chosen documents recovers queries from the documents they return
\cite{zhang2016all}. Sophos achieves forward privacy with a trapdoor permutation and a per-keyword
counter \cite{bost2016sophos}; Diana replaces the public-key operation with a constrained
pseudorandom function built from a GGM tree \cite{bost2017forward}. Backward privacy limits what
searches reveal about deleted entries and is classified into Types I to III
\cite{bost2017forward}; Ghareh Chamani et al. give constructions across this range
\cite{chamani2018new}. Amjad, Kamara, and Moataz generalize forward privacy to injection security
\cite{amjad2023injection}. These notions are defined over the transcript of operations observed
by a persistent adversary. In prior injection attacks the adversary supplies documents verbatim,
whereas in the setting of this paper the adversary controls only the task that causes the agent
to write.

\subsection{Snapshot Adversaries and History Independence}

A snapshot adversary observes the stored state at one or more points rather than the full
transcript. Amjad, Kamara, and Moataz define breach-resistant structured encryption, which
requires that a snapshot adversary learn less than a persistent one \cite{amjad2019breach}. Du,
Genkin, and Grubbs construct snapshot-oblivious RAMs \cite{du2022snapshot}, and Persiano and Yeo
establish lower bounds for such constructions \cite{persiano2023limits}. Grubbs, Ristenpart, and
Shmatikov observe that real database snapshots also contain logs, caches, and other artifacts that
leak past operations \cite{grubbs2017why}, and below the file system an SSD's flash translation
layer retains superseded pages after an in-place rewrite \cite{wei2011reliably}. History
independence, anticipated by Micciancio
\cite{micciancio1997oblivious} and defined by Naor and Teague \cite{naor2001anti}, requires that a
data structure's representation reveal nothing about the sequence of operations beyond its current
contents; Buchbinder and Petrank establish its inherent costs \cite{buchbinder2006history}. It has
been applied to file systems \cite{bajaj2013hifs} and combined with oblivious
storage and secure deletion \cite{roche2016oblivious}. Oblivious RAM
\cite{goldreich1996software,stefanov2013path} hides access patterns completely at a
polylogarithmic cost per access. The present work differs from breach resistance in that it
compares two cohorts of data under a single adversary, rather than two classes of adversary, and
it identifies the history-independence property that a dynamic encrypted vector index must
satisfy against that adversary.

\subsection{Private Vector Search}

Compass performs encrypted semantic search by running ORAM over an HNSW graph with client-side
traversal \cite{zhu2025compass}. Opal protects personal AI memory by confining data-dependent
reasoning to a trusted execution environment while an untrusted disk sees only fixed, oblivious
accesses; it handles continuous ingestion by piggybacking reindexing and capacity management on
each ORAM access \cite{kaviani2026opal}, so an oblivious store is not limited to static corpora.
What closes the derived-write channel there is oblivious access under a fixed budget, not the
enclave as such (\cref{prop:derived}). ONYX reduces the cost of
disk-oblivious nearest-neighbor search \cite{rathee2026onyx}. Pacmann \cite{zhou2025pacmann} and
Panther \cite{li2025panther} address private approximate nearest-neighbor search in client-server
settings. These systems protect against a compromised server or storage provider and, in Opal's
case, assume that the enclave is not compromised, an assumption that transient-execution attacks
have violated in practice \cite{vanbulck2018foreshadow}. None of them distinguishes records by the
time at which they were written relative to a compromise.

\subsection{Security of Agent Memory}

Most self-writing memories store records in plaintext, and the attacks of \cref{sec:intro}
either target integrity or, like MEXTRA \cite{mextra2025} and membership inference against
chat-agent memory \cite{mrmmia2026}, query the running agent, which access control addresses. A
recent survey confirms that the storage-layer confidentiality of a self-writing memory has not been
treated \cite{survey2026ltmsecurity}; this paper considers that adversary.
